\documentclass[11pt]{article}
\usepackage[margin=1in]{geometry}
\usepackage{amsmath,amssymb,amsthm,mathtools}
\usepackage{array}
\usepackage[hidelinks]{hyperref}

\newtheorem{theorem}{Theorem}[section]
\newtheorem{lemma}[theorem]{Lemma}
\newtheorem{proposition}[theorem]{Proposition}
\newcommand{\spE}{s^+}
\newcommand{\smE}{s^-}
\newcommand{\sig}{\sigma}
\newcommand{\one}{\mathbf 1}

\title{Positive Square Energy of Every Heptacyclic Cactus}
\author{}
\date{26 July 2026}

\begin{document}
\maketitle

\begin{abstract}
Let $G$ be a connected heptacyclic cactus of order $n$.  We prove
$s^+(G)>n$.  Sharp block-additive doubly nonnegative optimization reduces the
proof to the cycle multisets $T^6Q$ and $T^5PP$, where $T=C_3$, $P=C_5$, and
$Q$ is arbitrary.  A shared-triangle opening recurrence gives the quantitative
margins needed for both residuals.  If the shared-cut graph is disconnected,
$T^6Q$ is settled by an all-triangle leaf cluster.  For $T^5PP$, an exact
colored-partition audit has $46$ proper rows: $41$ are direct and five are
closed structurally, including an entry-sensitive degree dichotomy for
$T^5P\mid P$.  If all cycles form one shared-cut cluster, $T^6Q$ is settled by
opening a leaf $Q$ or splitting an internal $Q$, while $T^5PP$ is settled by
opening two pentagonal leaves or splitting an internal pentagon.  Independent
exact censuses reproduce the corrected totals $216,224,226,227,227$ for
$T^6Q$ and $560=557+3$ for $T^5PP$.  Every construction is a vertex partition
into induced subgraphs and permits arbitrary bridge connectors and arbitrary
trees attached at arbitrary vertices.
\end{abstract}

\section{Preliminaries and dependencies}

All graphs are finite, simple, and undirected.  A cactus is a connected graph
whose blocks are bridges or cycles.  Its cyclomatic number is the number of
cyclic blocks.  If the adjacency eigenvalues of $G$ are
$\lambda_1,\ldots,\lambda_n$, put
\[
 \spE(G)=\sum_{\lambda_j>0}\lambda_j^2,
 \qquad \smE(G)=\sum_{\lambda_j<0}\lambda_j^2,
 \qquad \sig(G)=\spE(G)-|V(G)|.
\]
For every vertex partition
$V(G)=V_1\mathbin{\dot\cup}\cdots\mathbin{\dot\cup}V_k$, induced-subgraph
superadditivity gives
\begin{equation}\label{eq:super}
 \spE(G)\geq\sum_{i=1}^k\spE(G[V_i]),
 \qquad
 \sig(G)\geq\sum_{i=1}^k\sig(G[V_i]).
\end{equation}
This is the theorem of Akbari, Kumar, Mohar, and Pragada~\cite{AKMP}.

Write $T=C_3$, $P=C_5$, and
\[
 \delta=\sec(\pi/5)-1=\sqrt5-2<\frac14,
 \qquad
 \delta_q=\sec(\pi/q)-1<1\quad(q\equiv1\pmod4).
\]
We use the following established packet bounds:
\begin{align}
 \sig(T)&>0,& \sig(P)&\geq-\delta,&
 \sig(C_q)&\geq-\delta_q,\label{eq:one}\\
 \sig(TT)&>1,& \sig(TP)&>1-\delta,&
 \sig(TC_q)&>1-\delta_q.\label{eq:two}
\end{align}
For the lower-rank classes,
\begin{align}
 \sig(H)&\geq0 &&\text{for every bicyclic or tricyclic cactus }H,
 \label{eq:small}\\
 \sig(H)&>0 &&\text{for every cactus of rank four, five, or six}.
 \label{eq:lower-rank}
\end{align}
For one shared-cut cluster we also use
\begin{align}
 \sig(PP)&\geq1-\frac4{3\sqrt{13}}>0,\label{eq:pp}\\
 \sig(TPP)&>6-2\sqrt5>\frac32.\label{eq:tpp}
\end{align}
Here a symbol records the cyclic blocks retained by the packet, not its tree
attachments.  Every estimate allows arbitrary finite trees attached at
arbitrary vertices; $TP$ and $TC_q$ also allow an arbitrary bridge connector.
The one-, two-, and three-cycle bounds are proved in the bicyclic and
tricyclic companion manuscripts~\cite{Bicyclic,Tricyclic}; the lower-rank
statements are the tetracyclic through hexacyclic cactus theorems
\cite{Tetracyclic,Pentacyclic,Hexacyclic}.  Bound \eqref{eq:pp} is the exact
shared-$PP$ estimate in~\cite{Bicyclic}, and \eqref{eq:tpp} is the shared
$C_3,C_5,C_5$ phase estimate in~\cite{Tricyclic}.

We also need the shared all-triangle estimates
\begin{equation}\label{eq:triangle-small}
 \sig(A_1)>0,
 \qquad \sig(A_2)>1,
 \qquad \sig(A_3)>2,
 \qquad \sig(A_4)>3,
\end{equation}
where the triangles of $A_r$ form one shared-cut cluster.  For $r\leq3$ these
follow from the favorable-cycle Sachs theorem~\cite{PackingTwo}.  For four
triangles the same theorem applies except to the central-triangle/three-petal
incidence; the exact grouped-Sachs matching injection in
\cite[Proposition~4.11]{PackingTwo} gives the same strict bound there.  Thus no
cycle-packing hypothesis is hidden in \eqref{eq:triangle-small}.

\section{Territories, intervals, and openings}

Join two cyclic blocks in the \emph{shared-cut graph} when they share a cut
vertex, and call its connected components shared-cut clusters.  In the
block-cut tree contract the incidence subtree of each cluster to one marked
node.  Every cluster, including a singleton cyclic block, is marked.  Take the
minimal subtree spanning the marked nodes and suppress only unmarked
degree-two nodes.  The result is the \emph{reduced cluster tree}.

\begin{lemma}[Arbitrary connector territories]\label{lem:territories}
Let $S_1,\ldots,S_k$ be pairwise vertex-disjoint connected subtrees of the
reduced cluster tree, and suppose every marked node belongs to exactly one
$S_i$.  Then $V(G)$ has a partition into connected induced territories
$G_1,\ldots,G_k$ such that $G_i$ retains exactly the cycles assigned to $S_i$.
Cuts occur on actual bridges.  Every connector remnant, Steiner branch, entry
vertex, and attached tree has exactly one owner.
\end{lemma}

\begin{proof}
Assign every remaining reduced-tree vertex to a nearest labelled subtree,
breaking ties by a fixed order.  Each label class in a tree is connected.
Every edge between classes expands in the block-cut tree to a path consisting
only of bridge blocks; cut one actual bridge on that path and assign its two
remnants to the adjacent owners.  At a Steiner branch assign the branch vertex
and one initial segment to one owner and cut an actual bridge on every other
incident segment.  A component outside the cyclic hull has a unique attachment
to it, since two attachments would create another cycle in the block-cut tree;
assign it whole to the owner of that attachment.  The resulting vertex sets
are disjoint, connected, and induced, and retain exactly the asserted cycles.
\end{proof}

For one shared-cut cluster, its \emph{cycle--cut incidence tree} is the
bipartite tree whose nodes are cyclic blocks and shared cyclic cuts.  A cycle
node is adjacent to the distinct cyclic cuts on that cycle.

\begin{lemma}[Intervals and private openings]\label{lem:intervals}
Let $C$ be a cycle node of a cycle--cut incidence tree.
\begin{enumerate}
\item If $C$ has $d\geq2$ distinct marked vertices, its vertices can be
partitioned into $d$ nonempty proper consecutive intervals, one owning each
mark.  Adjoining to an interval the corresponding component of $I-C$ gives a
connected induced territory.  The $d$ territories partition all vertices
assigned at $C$, every mark has one owner, and $C$ is retained by none.
\item If $v$ is private with respect to cyclic blocks, put $v$ and every
off-cycle tree branch rooted there into $F$.  Then $F$ is a nonempty tree and
$\sig(F)=-1$.  The path $C-v$ retains every cyclic cut of $C$ and belongs to
the complementary induced territory.
\end{enumerate}
Both constructions remain valid with arbitrary hanging trees.  An external
connector entering at a private cycle vertex may be assigned to either adjacent
interval; an entry at a marked cut or through an incidence branch follows the
unique interval owning that branch.
\end{lemma}

\begin{proof}
For the first construction, list the distinct marks cyclically and cut one
cycle edge in each selected gap between successive marks.  Assign unmarked
vertices in a gap to either adjacent interval.  Each interval is a proper path,
and each incidence branch follows its unique mark.  For the second, $C-v$ is a
path.  The vertex $v$ and all tree branches rooted at it form a nonempty tree.
For every nonempty tree $F$, bipartite spectral symmetry gives
$\spE(F)=|E(F)|=|V(F)|-1$, hence $\sig(F)=-1$.  Unique attachment of hanging
trees proves all ownership assertions.  Crossing edges are discarded only by
\eqref{eq:super}; no edge monotonicity is used.
\end{proof}

\section{The shared-triangle recurrence}

\begin{lemma}[Shared-triangle recurrence]\label{lem:triangle-recurrence}
Let $A_r$ be a connected cactus whose $r$ triangular cyclic blocks form one
shared-cut cluster, with arbitrary trees attached.  Then
\begin{equation}\label{eq:br}
 \sig(A_r)>b_r,
 \qquad
 \begin{array}{c|rrrrrrr}
 r&1&2&3&4&5&6&7\\ \hline
 b_r&0&1&2&3&2&1&0
 \end{array}.
\end{equation}
\end{lemma}

\begin{proof}
The values through four are \eqref{eq:triangle-small}.  For $r\geq5$, the
cycle--cut incidence tree has a leaf triangle.  It has one shared cyclic cut
and two vertices private with respect to cyclic blocks.  Open either private
vertex by Lemma~\ref{lem:intervals}.  The opened territory is a nonempty tree
of surplus $-1$.  Deleting the leaf cycle node, and suppressing its incident
cut if it becomes irrelevant, leaves the other $r-1$ triangle nodes connected
in one shared-cut incidence tree.  The complementary induced territory is
therefore an $A_{r-1}$ packet, with the surviving path fragment and all other
trees merely attached to it.  Consequently
\[
 \sig(A_r)\geq\sig(A_{r-1})-1.
\]
Starting from the strict bound $\sig(A_4)>3$ gives successively
$\sig(A_5)>2$, $\sig(A_6)>1$, and $\sig(A_7)>0$.
\end{proof}

The strict quantitative values $b_5=2$ and $b_6=1$ are essential below.
Qualitative lower-rank positivity is never used to pay a tree-opening cost or
a hostile odd-cycle deficit.

\section{The exact DNN residual classification}

For a connected edge-containing graph define
\[
 \kappa(G)=\sup_{\substack{M\succeq0,\ M\geq0\\\one^TM\one>0}}
 \frac{4\bigl(\sum_{uv\in E(G)}\sqrt{M_{uv}}\bigr)^2}
      {\one^TM\one}.
\]
The sharp cactus DNN theorem~\cite{DNN,LTZ} gives, for bridge-block count $b$
and cycle lengths $\ell_1,\ldots,\ell_r$,
\begin{equation}\label{eq:kappa}
 \kappa(G)=b+\sum_j\kappa(C_{\ell_j}),
 \qquad
 \kappa(C_\ell)=
 \begin{cases}
  \ell,&\ell\text{ even},\\
  \ell\sec^2(\pi/(2\ell)),&\ell\text{ odd},
 \end{cases}
 \qquad \smE(G)\leq\kappa(G).
\end{equation}
For completeness, fold each nonedge entry of a DNN matrix into its two
diagonal entries by adding $M_{uv}(e_u-e_v)(e_u-e_v)^T$.  Rotational averaging
on a cycle and its least Fourier eigenvalue give the cycle constants.
Orthogonal gluing of folded Gram systems at a one-vertex sum, followed by
Cauchy--Schwarz, proves block additivity.  Finally, writing $A=A_+-B$, the
Schur product theorem makes $B\circ B$ DNN and gives
\[
 \smE/2=-\sum_{uv\in E(G)}B_{uv}\leq\sum_{uv\in E(G)}|B_{uv}|,
 \qquad \one^T(B\circ B)\one=\smE;
\]
the definition of $\kappa$ yields $(\smE)^2\leq\kappa(G)\smE$.

Put
\[
 \epsilon_\ell=
 \begin{cases}
 0,&\ell\text{ even},\\
 \ell\tan^2(\pi/(2\ell)),&\ell\text{ odd}.
 \end{cases}
\]
The odd sequence is decreasing, since after putting $u=\pi/(2x)$ the relevant
derivative has the sign of $2u-\sin u\cos u>0$.  Moreover
\begin{equation}\label{eq:epsilon}
 \epsilon_3=1,
 \qquad \epsilon_5=5-2\sqrt5=:a,
 \qquad 3a<2,
 \qquad 2a>1,
 \qquad \epsilon_5+\epsilon_7<1.
\end{equation}
These comparisons are exact.  The middle two reduce, after squaring positive
sides, to $169<180$ and $80<81$.  Also
$\cos(\pi/7)>1-\tfrac12(22/49)^2>7/8$, whence
$\epsilon_7<7/15<2\sqrt5-4=1-\epsilon_5$; the last rational-radical
comparison squares to $4489<4500$.

For a heptacyclic cactus, block counting and \eqref{eq:kappa} give
\[
 b+\sum_{i=1}^7\ell_i=n+6,
 \qquad
 \sig(G)\geq6-\sum_{i=1}^7\epsilon_{\ell_i}.
 \tag{4.1}\label{eq:dnn}
\]
Indeed, $\spE+\smE=2|E(G)|=2n+12$ and
$\smE\leq n+6+\sum_i\epsilon_{\ell_i}$.  If at most four cycles are
triangles, the epsilon sum is at most $4+3a<6$.  With exactly five triangles,
an even remaining length contributes zero; if both remaining lengths are odd,
every pair except $5,5$ contributes at most
$\epsilon_5+\epsilon_7<1$.  The pair $5,5$ survives because $2a>1$.  With at
least six triangles the multiset has the form $\{3,3,3,3,3,3,q\}$, $q\geq3$.
Thus \eqref{eq:dnn} is strict outside exactly
\begin{equation}\label{eq:residual}
 T^6Q=\{3,3,3,3,3,3,q\}\quad(q\geq3),
 \qquad
 T^5PP=\{3,3,3,3,3,5,5\}.
\end{equation}

\section{Disconnected shared-cut graph: the $T^6Q$ family}

\begin{proposition}\label{prop:disconnected-t6q}
If the cycles are $T^6Q$ and the shared-cut graph is disconnected, then
$\spE(G)>|V(G)|$.
\end{proposition}

\begin{proof}
The reduced cluster tree is nontrivial and therefore has at least two marked
leaves.  At most one leaf cluster contains the designated block $Q$; this
remains true when $Q=T$, because that block remains designated.  Choose a
$Q$-free leaf cluster $A$.  It consists of $r$ triangles in one shared-cut
cluster, for some $1\leq r\leq6$.  Cut an actual bridge toward the rest and use
Lemma~\ref{lem:territories}.  The two connected induced territories are an
$A_r$ packet and a connected $(7-r)$-cyclic cactus containing $Q$.  The full
ledger is
\begin{center}
\begin{tabular}{c c c c}
$r$ & triangle leaf & complementary territory & total\\ \hline
$1$ & $>0$ & hexacyclic $>0$ & $>0$\\
$2$ & $>1$ & pentacyclic $>0$ & $>0$\\
$3$ & $>2$ & tetracyclic $>0$ & $>0$\\
$4$ & $>3$ & tricyclic $\geq0$ & $>0$\\
$5$ & $>2$ & bicyclic $\geq0$ & $>0$\\
$6$ & $>1$ & unicyclic $Q$ & $>0$
\end{tabular}
\end{center}
In the last row a nonhostile $Q$ is nonnegative.  If $q\equiv1\pmod4$, the
unicyclic territory, including all of its connector remnants and attached
trees, has surplus at least $-\delta_q$; hence the total is
$>1-\delta_q>0$.  This proves every reduced-tree topology and every connector
entry without an incidence census.
\end{proof}

\section{Disconnected shared-cut graph: all $46$ $T^5PP$ partitions}

Encode a shared-cut cluster by $(t,p)$, its numbers of triangles and pentagons.
Recursively choose the first nonzero pair, subtract it from $(5,2)$, and require
successive pairs to be lexicographically nondecreasing.  This gives exactly
$47$ colored set partitions, including the one-cluster partition, and hence
$46$ proper partitions.  Assign to each cluster the conservative exact-rational
ledger
\begin{center}
\begin{tabular}{c|cccccccccc}
packet&$T$&$P$&$TT$&$TP$&$PP$&$TTT$&$TPP$&$T^4$&$T^5$&$T^6$\\ \hline
bound&$>0$&$>-1/2$&$>1$&$>1/2$&$>0$&$>2$&$>3/2$&$>3$&$>2$&$>1$
\end{tabular}
\end{center}
and use $\geq0$ for every other rank-two or rank-three cluster and $>0$ for
every other cluster of rank four through six.  These are consequences of
\eqref{eq:one}--\eqref{eq:tpp}, \eqref{eq:lower-rank}, and
Lemma~\ref{lem:triangle-recurrence}.  A row is direct when the rational sum is
positive, or is zero with a strict summand.

The resulting $41$ direct rows, in canonical order, are
\begin{center}
\begin{minipage}{0.94\textwidth}\ttfamily\small
P|P|T|T|T|TT \quad P|P|T|T|TTT \quad P|P|T|TT|TT\par
P|P|T|TTTT \quad P|P|TT|TTT \quad P|P|TTTTT\par
P|T|T|T|T|TP \quad P|T|T|TP|TT \quad P|T|TP|TTT\par
P|T|TT|TTP \quad P|TP|TT|TT \quad P|TP|TTTT\par
P|TT|TTTP \quad P|TTP|TTT\par
PP|T|T|T|T|T \quad PP|T|T|T|TT \quad PP|T|T|TTT\par
PP|T|TT|TT \quad PP|T|TTTT \quad PP|TT|TTT \quad PP|TTTTT\par
T|T|T|T|TPP \quad T|T|T|TP|TP \quad T|T|T|TTPP\par
T|T|TP|TTP \quad T|T|TPP|TT \quad T|T|TTTPP\par
T|TP|TP|TT \quad T|TP|TTTP \quad T|TPP|TTT\par
T|TT|TTPP \quad T|TTP|TTP \quad T|TTTTPP\par
TP|TP|TTT \quad TP|TT|TTP \quad TP|TTTTP\par
TPP|TT|TT \quad TPP|TTTT \quad TT|TTTPP\par
TTP|TTTP \quad TTPP|TTT
\end{minipage}
\end{center}
Whole cluster territories and arbitrary connectors are supplied by
Lemma~\ref{lem:territories}.  Exact \texttt{Fraction} enumeration leaves the
following five rows, and no others:
\begin{equation}\label{eq:five-rows}
 \begin{gathered}
 P|P|T|T|T|T|T,\qquad
 TTP|P|T|T|T,\qquad
 TTTP|P|T|T,\\
 TTTTP|P|T,\qquad
 TTTTTP|P.
 \end{gathered}
\end{equation}

\begin{lemma}[The four singleton-triangle rows]\label{lem:four-rows}
The first four rows of \eqref{eq:five-rows} have positive surplus for every
reduced-tree topology and every connector entry.
\end{lemma}

\begin{proof}
Write these rows as
\[
 A_k\mid T\mid\cdots\mid T\mid P_1,
\]
where $A_0=P_0$, while for $k=2,3,4$ the packet
$A_k=T^kP_0$ is the unique nontrivial pentagon-containing cluster.  There is
at least one singleton triangle.  If one is a leaf of the reduced cluster tree,
cut the first actual bridge toward the rest.  The two induced territories are
a strict triangular unicyclic cactus and a strict connected hexacyclic cactus.

Suppose no singleton triangle is a leaf.  Every leaf of the minimal marked
tree is marked, so the only possible leaves are $A_k$ and $P_1$.  They are
therefore the two leaves and the reduced tree is their path.  Let $T_*$ be the
singleton triangle nearest $P_1$.  Cut an actual bridge immediately beyond the
terminal subtree containing $T_*$ and $P_1$.  Lemma~\ref{lem:territories}
gives
\[
 TP+(\text{connected pentacyclic cactus}).
\]
The first territory has surplus $>1-\delta>0$ and the second is strictly
positive.  This includes the all-singleton path
$P-T-T-T-T-T-P$.  The construction never splits $A_k$, so its incidence and
the entry by which the path reaches it are immaterial.
\end{proof}

\begin{lemma}[The $T^5P\mid P$ degree dichotomy]\label{lem:t5p-p}
Every cactus with cluster partition $T^5P_0\mid P_1$ has positive surplus.
\end{lemma}

\begin{proof}
Let $I$ be the incidence tree of the cluster $T^5P_0$ and put
$d=\deg_I(P_0)$.  Project the external connector entry to the first point where
it reaches the cyclic hull; an entry through a hanging branch follows the root
of that branch.

Suppose first that $d\geq2$.  The components of $I-P_0$ contain respectively
$r_1,\ldots,r_d\geq1$ triangles, with sum five.  Split $P_0$ into one proper
consecutive interval per component.  Adjoin the entire external connector and
$P_1$ to the interval owning its entry.  If that branch contains $r$ triangles,
its retained territory has type $T^rP_1$: for $r=1$ it is a strict $TP$
packet, for $r=2$ it is a nonnegative tricyclic cactus, and for $r\geq3$ it is
strictly positive by \eqref{eq:lower-rank}.  Every other interval territory is
an $A_{r_j}$ packet and is strict.  Since $d\geq2$, at least one such other
territory exists, so the sum is strictly positive even when the entry territory
has only the nonnegative tricyclic bound.  If the entry is a private vertex of
$P_0$, assign it and the connector to either adjacent interval; at a shared cut
or through a triangular branch the incidence component determines its unique
owner.  This also covers the saturated case $d=5$.

Now suppose $d=1$.  Choose a private vertex $v_0$ of $P_0$ away from the
connector entry and a private vertex $v_1$ of $P_1$ away from its entry.  Such
choices exist: $P_0$ has four vertices private with respect to cyclic blocks,
and $P_1$ has no shared cyclic cut.  Open both pentagons at these vertices.
The two opened territories $F_0,F_1$ are disjoint nonempty trees with surplus
$-1$ each.  All remaining vertices form a connected induced territory $H$:
the path $P_0-v_0$ remains attached at its unique cyclic cut, the path
$P_1-v_1$ remains joined through the connector, and the choices away from the
entries preserve the connector route.  The only cycles of $H$ are the five
triangles.  Deleting the incidence-leaf $P_0$ leaves them in one shared-cut
cluster, so Lemma~\ref{lem:triangle-recurrence} gives
\[
 \sig(G)\geq\sig(H)+\sig(F_0)+\sig(F_1)>2-1-1=0.
\]
Thus both possible degrees, every cyclic order, and every connector entry are
covered.
\end{proof}

\begin{proposition}\label{prop:disconnected-t5pp}
If the cycles are $T^5PP$ and the shared-cut graph is disconnected, then
$\spE(G)>|V(G)|$.
\end{proposition}

\begin{proof}
The colored recursion gives all $46$ proper partitions.  The packet ledger
settles the displayed $41$ direct rows.  Lemma~\ref{lem:four-rows} settles the
first four remaining rows and Lemma~\ref{lem:t5p-p} settles the fifth.  Hence
every colored partition, reduced-tree topology, connector entry, and internal
cluster incidence is covered.
\end{proof}

\section{One fully shared $T^6Q$ cluster}

\begin{proposition}[Leaf $Q$ or internal $Q$ split]\label{prop:shared-t6q}
If all seven $T^6Q$ cycles form one shared-cut cluster, then
$\spE(G)>|V(G)|$.
\end{proposition}

\begin{proof}
Let $I$ be the cycle--cut incidence tree and inspect the designated node $Q$.
If $\deg_I(Q)=1$, open a vertex of $Q$ private with respect to cyclic blocks.
Such a vertex exists because $q\geq3$.  The opened territory is a nonempty
tree of surplus $-1$.  Deleting the incidence-leaf $Q$ leaves the six triangle
nodes connected in one shared-cut cluster, so the complementary induced
territory is an $A_6$ packet.  Therefore
\[
 \sig(G)\geq\sig(A_6)-1>1-1=0.
\]

If $d=\deg_I(Q)\geq2$, split $Q$ into one proper consecutive interval per
component of $I-Q$.  Each component contains at least one triangle: the cut
neighbor of $Q$ has degree at least two, and no other nontriangle exists.  If
their triangle counts are $r_1,\ldots,r_d$, then $r_j\geq1$ and
$\sum_jr_j=6$.  The exact induced territories are $A_{r_1},\ldots,A_{r_d}$,
with the path intervals of $Q$ and all hanging trees absorbed as attachments.
By Lemma~\ref{lem:triangle-recurrence},
\[
 \sig(G)\geq\sum_j\sig(A_{r_j})>\sum_j b_{r_j}\geq0.
\]
Strictness follows from the strict estimate for every branch.  This case
includes a saturated $Q$ hub and all dispersed branch patterns.  The
seven-cycle bouquet lies in the first case: opening $Q$ leaves the six-triangle
common-cut packet and closes it with one tree cost.
\end{proof}

For comparison and as an independent finite certificate, let $c$ be the
number of shared cyclic cuts.  The incidence tree has $c+6$ edges and
\[
 \sum_{x\text{ cut}}(\deg_I(x)-1)=6,
 \qquad 1\leq c\leq6.
\]
Quotienting by permutations of the six triangles and cut nodes, while keeping
$Q$ distinguished even when $q=3$, gives the corrected census
\begin{center}
\begin{tabular}{c|rrrrrr|r}
$Q$ capacity&$c=1$&$c=2$&$c=3$&$c=4$&$c=5$&$c=6$&total\\ \hline
$q=3$&1&8&33&71&74&29&216\\
$q=4$&1&8&33&73&77&32&224\\
$q=5$&1&8&33&73&78&33&226\\
$q=6$&1&8&33&73&78&34&227\\
$q\geq7$&1&8&33&73&78&34&227
\end{tabular}
\end{center}
An exact conservative packet test finds a safe ordinary one-cycle split in
every nonbouquet class, respectively $215,223,225,226,226$ classes.  The sole
ordinary-split exception in every regime is the seven-cycle bouquet already
handled in Proposition~\ref{prop:shared-t6q}.  These counts are a certificate
of the structural proof, not a premise of it.

\section{One fully shared $T^5PP$ cluster}

\begin{proposition}[Two pentagonal leaves or an internal pentagon]
\label{prop:shared-t5pp}
If all seven $T^5PP$ cycles form one shared-cut cluster, then
$\spE(G)>|V(G)|$.
\end{proposition}

\begin{proof}
Let $I$ be the incidence tree and first suppose both pentagon nodes are leaves.
Open one private vertex on each pentagon.  The opened territories are disjoint
nonempty trees of surplus $-1$.  Successively deleting the two leaf cycle nodes
from the tree leaves the five triangle nodes connected in one shared-cut
incidence tree.  The complementary territory is therefore an $A_5$ packet,
and
\[
 \sig(G)\geq\sig(A_5)-2>2-2=0.
\]
The two pentagons are distinct blocks, so their opened territories are
disjoint; each private choice can include all tree branches rooted there.

Otherwise an internal pentagon $P_0$ exists.  Put
$2\leq d=\deg_I(P_0)\leq5$ and split $P_0$ into one proper consecutive
interval per component of $I-P_0$.  There is a unique branch $B$ containing
the other pentagon $P_1$; let $a$ be the number of triangles in $B$.

If $a\geq1$, then $a\leq4$, because at least one other incidence component
contains a cycle and every other cycle is triangular.  The $B$ territory has
type $T^aP$: it is strictly positive for $a=1$ by \eqref{eq:two}, nonnegative
for $a=2$ by \eqref{eq:small}, and strictly positive for $a=3,4$ by
\eqref{eq:lower-rank}.  At least one other branch is an all-triangle packet and
is strict by Lemma~\ref{lem:triangle-recurrence}.  The total is positive.

If $a=0$, then $B$ is a singleton pentagon and contributes at least $-\delta$.
The five triangles occupy at most $d-1\leq4$ other nonempty branches.  Some
branch therefore contains at least two triangles; its shared-triangle margin is
$>1$ by Lemma~\ref{lem:triangle-recurrence}.  Every remaining triangular
branch is strict, and the total is $>1-\delta>0$.  This includes a saturated
pentagon hub: if the other pentagon is a singleton petal, the five triangles in
at most four other branches force the required $TT$-or-larger packet.  The two
cases exhaust the incidence tree and every interval has one owner by
Lemma~\ref{lem:intervals}.
\end{proof}

The independent exact census provides a second completeness check.  Quotient
by $S_5\times S_2\times S_c$ and test every cycle deletion against the retained
incidence hypotheses of the packet ledger.  The exact result is
\begin{center}
\begin{tabular}{c|rrrrrr|r}
$c$&1&2&3&4&5&6&total\\ \hline
colored incidence trees&1&12&68&177&211&91&560\\
SAFE ordinary split&0&11&67&177&211&91&557
\end{tabular}
\end{center}
The three canonical ordinary-split exceptions are exactly
\begin{align}
 &((0,7),(1,7),(2,7),(3,7),(4,7),(5,7),(6,7)),
 \label{eq:u1}\\
 &((0,7),(1,7),(2,7),(3,7),(4,7),(5,7),(0,8),(6,8)),
 \label{eq:u2}\\
 &((0,7),(1,7),(2,7),(3,7),(4,7),(0,8),(5,8),(1,9),(6,9)),
 \label{eq:u3}
\end{align}
where $0,\ldots,4=T$, $5,6=P$, and cut nodes begin at $7$.  They are the
seven-cycle bouquet, a common-cut $T^5P$ core with a $TP$ tail, and a
five-triangle common-cut core with two separate pentagon tails.  In all three,
both pentagons are incidence leaves.  Thus the first case of
Proposition~\ref{prop:shared-t5pp} opens them and leaves one concentrated
$A_5$ packet.  In particular, the exact $560=557+3$ certificate has no
unresolved saturated hub.  The structural proposition proves all classes
without relying on the census.

\section{Main theorem}

\begin{theorem}\label{thm:main}
Every connected heptacyclic cactus $G$ on $n$ vertices satisfies
\[
 \boxed{\spE(G)>n}.
\]
\end{theorem}

\begin{proof}
The exact DNN estimate \eqref{eq:dnn} is strict unless the cycle multiset is
$T^6Q$ or $T^5PP$.  For $T^6Q$, a disconnected shared-cut graph is handled by
Proposition~\ref{prop:disconnected-t6q}, while one shared-cut cluster is handled
by Proposition~\ref{prop:shared-t6q}.  For $T^5PP$, the corresponding cases are
Propositions~\ref{prop:disconnected-t5pp} and
\ref{prop:shared-t5pp}.  These alternatives exhaust every cycle multiset and
every shared-cut incidence.

Lemmas~\ref{lem:territories} and \ref{lem:intervals} assign every actual
connector segment, Steiner branch, entry vertex, shared cyclic cut, and hanging
tree to exactly one territory.  Private openings are paid only by the explicit
shared-triangle margins of Lemma~\ref{lem:triangle-recurrence}.  Every spectral
comparison is applied to a vertex partition into induced subgraphs; no edge
monotonicity, tree relocation, bounded connector length, or special attachment
position is assumed.
\end{proof}

Since a connected heptacyclic graph has $|E(G)|=n+6$, this proves the strict
conclusion of the Akbari--Kumar--Mohar--Pragada--Zhang conjecture
\cite[Conjecture~1.2]{AKMPZ} for every heptacyclic cactus.

\appendix
\section{Exact reproduction and proof objects}\label{app:reproduction}

From the repository root run, using only the Python standard library,
\begin{verbatim}
python research/heptacyclic-t5pp-disconnected-partition-audit.py
python research/heptacyclic-t6q-incidence-census.py
python research/heptacyclic-tttttpp-incidence-census.py
\end{verbatim}
All numerical proof comparisons in these scripts use
\texttt{fractions.Fraction}; no floating-point comparison is a premise.  The
partition script asserts
\begin{verbatim}
colored partitions including one cluster: 47
proper colored partitions: 46
direct packet rows: 41
topology/entry rows: 5
\end{verbatim}
and asserts exactly the five rows in \eqref{eq:five-rows}.

The $T^6Q$ script asserts, by cut count $c=1,\ldots,6$,
\begin{verbatim}
q=3:  1, 8, 33, 71, 74, 29 = 216
q=4:  1, 8, 33, 73, 77, 32 = 224
q=5:  1, 8, 33, 73, 78, 33 = 226
q=6:  1, 8, 33, 73, 78, 34 = 227
q>=7: 1, 8, 33, 73, 78, 34 = 227
\end{verbatim}
and safe-split totals $215,223,225,226,226$, with the bouquet as the unique
exception in each regime.  Generation is by exhaustive leaf extension of the
rank-six $T^5Q$ incidence classes, followed by an independent center-rooted
color-preserving tree code.

The $T^5PP$ script asserts
\begin{verbatim}
colored trees by cut count:       1, 12, 68, 177, 211, 91 = 560
SAFE-resolved trees by cut count: 0, 11, 67, 177, 211, 91 = 557
unresolved structural classes: bouquet, T^5P core with TP tail,
                               T^5 core with two P tails
\end{verbatim}
It checks all $7\cdot560=3920$ cycle choices and verifies the retained
shared-cut hypotheses before using concentrated packet bounds.  Its rational
proxies include $P>-1/4$, $TP>3/4$, shared $TTP>7/4$, and shared
$TPP>3/2$.  These follow exactly from $\delta<1/4$ and
$6-2\sqrt5>3/2$, certified by squaring positive rational bounds.

The scripts enumerate finite colored incidence objects or colored cluster
partitions.  They do not infer inducedness from abstract trees.  Induced
realization, cyclic intervals, connector ownership, arbitrary attached trees,
and the structural repairs are proved in the body of the manuscript.  The
current source proof objects are
\begin{verbatim}
research/heptacyclic-dnn-residuals-2026-07-26.md
research/heptacyclic-t6q-disconnected-2026-07-26.md
research/heptacyclic-t5pp-disconnected-complete-audit-2026-07-26.md
research/heptacyclic-residual-sacrifice-splitting-lemmas-2026-07-26.md
research/heptacyclic-cactus-complete-synthesis-2026-07-26.md
research/heptacyclic-t5pp-disconnected-partition-audit.py
research/heptacyclic-t6q-incidence-census.py
research/heptacyclic-tttttpp-incidence-census.py
\end{verbatim}
The corrected totals displayed here are generated by these current scripts;
no earlier pre-fix census count is used.

\begin{thebibliography}{99}
\bibitem{AKMP}
S.~Akbari, H.~Kumar, B.~Mohar, and S.~Pragada,
\emph{A linear lower bound for the square energy of graphs},
Electron. J. Combin. 32(3) (2025), Paper No. P3.53.

\bibitem{AKMPZ}
S.~Akbari, H.~Kumar, B.~Mohar, S.~Pragada, and S.~Zhang,
\emph{Refinement of a conjecture on positive square energy of graphs},
arXiv:2506.07264, 2025.

\bibitem{Bicyclic}
Companion manuscript, \emph{Positive Square Energy of Every Bicyclic Cactus},
2026; repository path \texttt{all-bicyclic-cacti/paper.tex}.

\bibitem{Tricyclic}
Companion manuscript, \emph{Positive Square Energy of Every Tricyclic Cactus},
2026; repository path \texttt{all-tricyclic-cacti/paper.tex}.

\bibitem{Tetracyclic}
Companion manuscript, \emph{Positive Square Energy of Every Tetracyclic
Cactus}, 2026; repository path \texttt{all-tetracyclic-cacti/paper.tex}.

\bibitem{Pentacyclic}
Companion manuscript, \emph{Positive Square Energy of Every Pentacyclic
Cactus}, 2026; repository path \texttt{all-pentacyclic-cacti/paper.tex}.

\bibitem{Hexacyclic}
Companion manuscript, \emph{Positive Square Energy of Every Hexacyclic
Cactus}, 2026; repository path \texttt{all-hexacyclic-cacti/paper.tex}.

\bibitem{DNN}
Companion manuscript, \emph{The Optimized Liu--Tang--Zhang Constant for Cactus
Graphs}, 2026; repository path \texttt{sharp-cactus-dnn/paper.tex}.

\bibitem{PackingTwo}
Companion manuscript, \emph{Square Energy from Cycle Packing and Block-Graph
Territories}, 2026; repository path
\texttt{packing-two-square-energy/paper.tex}.

\bibitem{LTZ}
Y.~Liu, Q.~Tang, and S.~Zhang,
\emph{The positive and negative square-energy conjecture},
arXiv:2607.18031, 2026.
\end{thebibliography}
\end{document}
